% =====================================================================
% SDUBeamer 完整可编译示例
% ---------------------------------------------------------------------
% 编译方式（推荐，支持中文）：
%   cd examples && TEXINPUTS=..:$TEXINPUTS xelatex demo.tex
%   TEXINPUTS=..:$TEXINPUTS xelatex demo.tex   # 建议执行两遍以解决交叉引用
%
% 亦可从仓库根目录编译：
%   xelatex -output-directory=examples examples/demo.tex
%   （模板 sdubeamer.sty 位于仓库根目录，需确保可被 TEXINPUTS 找到）
%
% 亦可用 latexmk：
%   latexmk -xelatex -cd examples/demo.tex
%
% 说明：
%   - 本示例覆盖了标题页、章节页、内容帧、分栏、列表、
%     数学公式、代码块、引用、表格等常见 Beamer 元素。
%   - 模板默认开启中文支持（基于 ctex + fandol 字体集），
%     因此必须使用 xelatex 编译。
% =====================================================================
\documentclass[aspectratio=169]{beamer}

% 加载 SDUBeamer 模板，可选主题色 blue(默认)/red/cyan/green/purple
\usepackage[color=blue]{sdubeamer}

% 如需自定义字体集（例如 Windows 系统自带字体），可覆盖 fontset 选项：
% \usepackage[fontset=windows]{sdubeamer}

% 若为纯英文场景，可关闭中文支持：
% \usepackage[chinese=false]{sdubeamer}

% 代码高亮（可选）
\usepackage{listings}
\lstset{
  basicstyle=\ttfamily\footnotesize,
  breaklines=true,
  frame=single,
  numbers=left,
  numberstyle=\tiny\color{gray},
  keywordstyle=\color{sdu@primary}\bfseries,
}

% =====================================================================
\begin{document}

% --------------------------- 标题页 ---------------------------
\title{SDUBeamer 幻灯片模板}
\subtitle{山东大学 Beamer 主题使用示例}
\author{张三}
\institute{山东大学 计算机科学与技术学院}
\date{2025 年 6 月}

\begin{frame}
  \titlepage
\end{frame}

% --------------------------- 目录 ---------------------------
\begin{frame}{目录}
  \tableofcontents
\end{frame}

% =====================================================================
\section{模板介绍}

\begin{frame}
  \sectionpage
\end{frame}

\begin{frame}{SDUBeamer 简介}
  \begin{itemize}
    \item 面向山东大学学位论文答辩的 Beamer 主题
    \item 支持本科生、硕士研究生、博士研究生答辩场景
    \item 内置中文支持，基于 \texttt{ctex} 宏包
  \end{itemize}

  \vspace{0.5cm}
  \begin{block}{主要特性}
    \begin{enumerate}
      \item 多种主题色可选（\texttt{blue}/\texttt{red}/\texttt{cyan}/\texttt{green}/\texttt{purple}）
      \item 免安装的 \texttt{fandol} 字体集，跨平台可编译
      \item 内置标题页、章节页等自定义模板
    \end{enumerate}
  \end{block}
\end{frame}

% =====================================================================
\section{功能演示}

\begin{frame}
  \sectionpage
\end{frame}

\begin{frame}{分栏布局}
  \begin{columns}
    \begin{column}{0.48\textwidth}
      \begin{block}{左侧栏}
        使用 \texttt{columns} 环境可以实现左右分栏布局，
        适合图文混排或对比展示。
      \end{block}
      \begin{alertblock}{警示框}
        使用 \texttt{alertblock} 展示需要强调的内容。
      \end{alertblock}
    \end{column}
    \begin{column}{0.48\textwidth}
      \begin{exampleblock}{示例框}
        使用 \texttt{exampleblock} 展示示例或结论。
      \end{exampleblock}
      \begin{itemize}
        \item 要点一
        \item 要点二
        \item 要点三
      \end{itemize}
    \end{column}
  \end{columns}
\end{frame}

\begin{frame}{数学公式}
  高斯分布的概率密度函数为：
  \[
    f(x) = \frac{1}{\sqrt{2\pi\sigma^2}}
    \exp\left(-\frac{(x-\mu)^2}{2\sigma^2}\right)
  \]

  欧拉恒等式：
  \[
    e^{i\pi} + 1 = 0
  \]
\end{frame}

\begin{frame}[fragile]{代码展示}
  \begin{lstlisting}[language=C++]
#include <iostream>

int main() {
    std::cout << "Hello, SDUBeamer!" << std::endl;
    return 0;
}
  \end{lstlisting}
\end{frame}

\begin{frame}{表格}
  \begin{table}
    \centering
    \caption{示例数据}
    \begin{tabular}{lcc}
      \hline
      项目 & 数量 & 备注 \\
      \hline
      章节 & 3     & 结构清晰 \\
      帧数 & 8     & 演示充分 \\
      代码块 & 1   & 语法高亮 \\
      \hline
    \end{tabular}
  \end{table}
\end{frame}

\begin{frame}{逐条显示}
  \begin{itemize}[<+->]
    \item 第一条要点
    \item 第二条要点
    \item 第三条要点
  \end{itemize}
\end{frame}

% =====================================================================
\section{总结}

\begin{frame}
  \sectionpage
\end{frame}

\begin{frame}{总结}
  \begin{block}{结语}
    本示例演示了 SDUBeamer 模板的主要使用方式。
    更多细节请参考项目的 \texttt{README.md} 文档。
  \end{block}
  \begin{center}
    \large 感谢聆听！\\
    \normalsize 欢迎提问与交流。
  \end{center}
\end{frame}

\end{document}
